\section{朝廷之外的权威：边防、经学与社会网络}
\label{sec:eastern-han-court-society}

明、章时期的守成并不表示雒阳拥有取之不尽的资源。西域驻军、羌地冲突、北边防务和宫廷人事都要求同一批粮食、将领和可信中介。窦宪北出、班超经营西域与甘英西使，显示东汉仍可把军队和使节送到远处；燕然山铭的功绩语言、甘英所闻海路与“大秦”消息，却不能直接换成歼敌数、精确路线或直接外交的事实。

和帝借宦官清除窦宪，解决了外戚把持宫门的即时风险，也让守门、传诏和宫廷安全成为宦官更可进入的权力通道。这不是“好人”与“坏人”的替换：幼帝、后族、近臣和官僚都要解决谁能保护皇帝、传递命令和任用将领。临时解决改变下一场危机的可用关系，外戚与宦官反复交替，正是这些高风险接口的政治后果。

\subsection{谁能规定可公开的知识}

白虎观讨论、太学与熹平石经将经义、书写和官学权威置于可见的位置。它们表明朝廷仍试图规定哪些文字、仪式与师说可以公开确认，不证明全国读书人都接受一套课程。石经残石、传本和遗址各自校验不同层面的事实，不能替我们画出所有士人的社会网络。鸿都门学纳入辞赋、书画与尺牍之士，也说明声望并不只在“朝廷儒学”一条路上产生。

\eventanchor{EH-0079-WHITE-TIGER-HALL}{建初四年白虎观讨论}和\eventanchor{EH-0175-XIPING-STONE-CLASSICS}{熹平石经的刊立}，把经义从师承与口授的一部分，变成朝廷能够召集、校定和陈列的公共权威。《白虎通》的编定同现场记录的关系不清，残石也不能告诉我们所有学人的读法；不过，两种材料合起来足以说明文字、碑石、太学空间和官署判断正在彼此借力。被刻定的知识为士人提供了接近官学的道路，也为不被承认的师说留下新的边界。

楚王英案中黄老与佛教供养的字样，是佛教进入官方文字的早期线索，不是全国已有统一佛教制度的证据。师友、家族、书写技术、地方评价与宗教结社能够形成跨县界的声望，许多网络同时与州郡、豪强相接，并不简单属于“反国家”。正史保留宫廷语言最完整，碑刻、墓葬和宗教文本让我们看见的则是不同的物质与社会尺度。

\eventanchor{EH-0065-CHU-PRINCE-YING}{楚王英案}所存诏语中的黄老、佛教供养用字，是一条早期而狭窄的线索：处分和徙置可核，术语的文本层次、信众规模和具体实践范围却仍有争议。与其把它写成某种宗教已经覆盖帝国，不如从一个王国的指控、财物和书写看见宫廷如何命名可疑的结社与供养。宗教活动、家族资源和地方声望能够互相借用，却不能仅凭官方案件就还原为同一种组织。

一五九年梁冀集团被诛、一六六年党锢开始，窦武与陈蕃又在一六八年失败。党人的名单、地域与持续时间并不完全一致，不能写成铁板一块的阵营；逮捕和禁锢的事实说明朝廷日益把本可提供地方评价的士人网络视为竞争资源。边地供给、宫门保护和公共声望的争夺至此已经难以彼此隔开。

\begin{textwindow}[权威并非只在宫门内]
《后汉书》的诏令、列传与制度志让朝廷的声音最为响亮；石经残片、碑刻、墓葬和早期宗教文本不能替代正史，却使读者看见谁在书写、捐献、迁徙和组织互助。材料不足以复原完整社会网络，正因如此，不能让宫廷斗争吞没社会史。
\end{textwindow}

\begin{sourcenote}[功绩与知识的限度]
《后汉书》帝纪、窦宪传、班超传、宦者列传、儒林传和党锢列传为本节的年序骨架；燕然铭、石经残石和宗教材料提供局部交叉。功绩语、传闻、党人名单与人物动机均受文本位置限制，不作超过材料承载力的推断。
\end{sourcenote}
\subsection{账本增补的编年桥}
这些首次出现的锚点以账本的纪年、行动者与史料限度串联；它们只承载可核的桥接，不把简略记载延伸为未经证实的场景。

\noindent 建武中元二年（57年），东汉面对建武时期主要竞争政权已被压制，刘庄作为长期确立的太子拥有相对稳定的继承基础；\eventanchor{EH-0057-MINGDI-ACCESSION}{光武帝去世，太子刘庄即位，是为明帝}，其后东汉由开国军事动员转向守成朝廷，明帝朝的礼制、边疆和官僚政策在既有框架内调整。记事依据《后汉书·光武帝纪下》；《后汉书·明帝纪》；《后汉纪》卷二十八；继位年月可靠；光武晚年继承安排的细节有争议。
\par

\noindent 永平十八年（75年）的东汉发生\eventanchor{EH-0075-ZHANGDI-ACCESSION}{明帝去世，太子刘炟即位，是为章帝}。明帝在位时皇位继承已较明确，朝廷需要在边疆危机中保持中枢任命和财政连续构成其条件，章帝即位后由窦氏等外戚影响朝政，东汉宫廷权力结构进入新的调整期随之进入后续年序。《后汉书·明帝纪》；《后汉书·章帝纪》；《后汉纪》卷三十六提供这一线索，但继位年月可靠；明帝遗命与辅政安排的细节有异文。
\par

\noindent 在章和二年（88年），\eventanchor{EH-0088-HEDI-ACCESSION}{章帝去世，皇子刘肇即位，是为和帝，窦太后临朝}把东汉置于新帝年少，窦氏凭皇太后身份和禁中网络取得辅政位置与窦氏成为中枢人事与军事任命的核心，后续北伐功臣和外戚集团的关系更趋紧张之间。可据《后汉书·章帝纪》；《后汉书·和帝纪》；《后汉纪》卷四十九追踪其大体经过；继位事实可靠；章帝后期储位变动和窦氏权力配置有争议。
\par

\noindent 东汉在永元四年（92年）的\eventanchor{EH-0092-DOU-CLAN-PURGE}{和帝联同宦官郑众诛除窦宪集团，结束窦太后家族的军政主导}，并非孤立一事：窦宪以北伐和禁军权势压过皇帝，和帝需借近侍和宫廷通道打破外戚控制，且和帝得以亲政，宦官因参与政变进入更关键的政治位置，外戚—宦官竞争的结构加深。材料见《后汉书·和帝纪》；《后汉书·窦宪传》；《后汉书·宦者列传》；政变事实可靠；和帝自主程度与郑众角色有争议。
\par

\noindent \eventanchor{EH-0105-HEDI-DEATH}{和帝去世，襁褓中的刘隆即位，是为殇帝，邓太后临朝}发生于元兴元年（105年），行动者为东汉。它回应和帝早逝且继承人年幼，邓氏以皇太后身份掌握诏令、禁军和人事渠道，并把东汉再次进入外戚摄政，皇帝年龄使宫廷派系和地方危机更难由常规君主权威调节留给随后的人群和官署。《后汉书·和帝纪》；《后汉书·殇帝纪》；《后汉书·邓皇后纪》可作核验，惟继位年月可靠；和帝病逝与继嗣选择的宫廷过程有争议。
\par

\noindent 从婴儿殇帝夭亡，朝廷需在刘氏宗室中选择可被邓氏与大臣接受的继承人出发，东汉于延平元年（106年）出现\eventanchor{EH-0106-AN-DI-ACCESSION}{殇帝去世，清河王刘祜即位，是为安帝，邓太后继续临朝}；安帝即位维持皇统连续，却使邓氏摄政延长，外戚权力与皇帝成年后的关系埋下冲突。这里采用《后汉书·殇帝纪》；《后汉书·安帝纪》；《后汉书·邓皇后纪》的记载范围，不越过继统事实可靠；废立候选与邓氏决策过程有争议。
\par

\noindent 元初三年（116年），东汉的\eventanchor{EH-0116-DENG-SUI-DEATH}{邓太后去世，安帝开始亲政并贬逐部分邓氏成员}使安帝成年后长期受邓氏限制，太后去世移除了维持摄政平衡的核心人物转为可见的行动，皇帝重新控制人事，邓氏网络被削弱；外戚摄政的退出再次以报复性重组而非制度化交接收场。其依据为《后汉书·安帝纪》；《后汉书·邓皇后纪》；《后汉纪》卷七十七，太后卒年与政局转换可靠；安帝报复的范围和个案罪名有争议。
\par

\noindent \eventanchor{EH-0125-AN-DI-SUCCESSION}{安帝去世，阎后先立北乡侯，宦官孙程等旋迎立济阴王刘保为顺帝}这一延光四年（125年）的节点，关联东汉、安帝遗嗣与阎后集团的利益冲突使继承程序分裂，阎氏缺少足以压制宫中近侍的联盟和顺帝即位，宦官以拥立功获得封赏和政治资本，宫廷继承更依赖禁中武力与信息通道。《后汉书·安帝纪》；《后汉书·顺帝纪》；《后汉书·孙程传》记录了它的基本轮廓；继位转折可靠；北乡侯死因和孙程等人行动细节有争议。
\par

\noindent 东汉在建康元年（144年）处理顺帝留下年幼继承人，梁氏凭皇后、外戚和朝廷官职控制继承安排时，出现\eventanchor{EH-0144-SHUNDI-DEATH}{顺帝去世，幼帝刘炳即位，是为冲帝，梁太后与梁冀临朝}；梁太后和梁冀进入权力中心，外戚对禁军和任官的影响显著扩大。可核材料为《后汉书·顺帝纪》；《后汉书·冲帝纪》；《后汉书·梁冀传》，且继位年月可靠；梁氏在军政中的分工有争议。
\par

\noindent 永嘉元年（145年）的\eventanchor{EH-0145-CHONGDI-DEATH}{冲帝夭亡，梁太后立八岁的刘缵为质帝，梁冀继续掌政}，让东汉的冲帝在位短暂，梁氏必须迅速另择宗室幼主以保持摄政结构在具体事务中显形；连续的幼帝继承使皇帝权威进一步空洞化，梁冀的实际权力和朝臣依附关系更加强化。《后汉书·冲帝纪》；《后汉书·质帝纪》；《后汉书·梁冀传》所能支持的范围仍受冲帝卒与质帝即位可靠；选立动机与宫廷谈判有争议限制。
\par

\noindent 对东汉而言，\eventanchor{EH-0146-ZHIDI-DEATH-HUANDI}{质帝被毒杀后，梁太后立刘志为桓帝，梁冀继续控制朝政}发生在本初元年（146年）：质帝对梁冀的公开批评威胁摄政集团，梁氏需寻找较易控制的成年宗室，随后桓帝登位却受梁冀掣肘，毒杀指控成为理解外戚专权的重要叙事，亦需区分史家谴责与可证事实。参看《后汉书·质帝纪》；《后汉书·桓帝纪》；《后汉书·梁冀传》；废立与桓帝即位可靠；毒杀责任主要见于后出史家叙事，证据有争议。
\par

\noindent 李固等主张限制梁氏并批评连续幼帝政治，触及梁冀对人事和继统的控制之下，东汉于建和元年（147年）经历\eventanchor{EH-0147-LI-GU-PURGE}{梁冀诬杀李固、杜乔等反对者，朝廷清流官僚受挫}，并面对反梁官僚被清除，士人以名节评价政治的语言更趋尖锐，朝廷监督机制进一步失衡。此处依《后汉书·李固传》；《后汉书·杜乔传》；《后汉纪》卷八十三，处置事实可靠；党派标签和密谋情节有争议。
\par

\noindent 延熹二年（159年），东汉面对梁冀长期控制朝政并压制皇帝，桓帝只能借禁中宦官发动突袭；\eventanchor{EH-0159-LIANG-JI-FALL}{桓帝与单超等宦官诛灭梁冀集团，梁氏外戚专权结束}，其后梁氏被清算，宦官因功受封掌权；外戚问题未被制度解决，而转化为宦官—士人冲突。记事依据《后汉书·桓帝纪》；《后汉书·梁冀传》；《后汉书·宦者列传》；政变年月可靠；桓帝主导程度与查抄财产数额有争议。
\par

\noindent 延熹九年（166年）的东汉发生\eventanchor{EH-0166-FIRST-PARTISAN-PROHIBITION}{桓帝朝以李膺等结党为名逮捕太学生和官僚，形成第一次党锢}。宦官与依附官僚担忧太学生、名士的舆论和荐举网络挑战其任官利益构成其条件，清议士人与宦官集团的对立公开化；“党人”是国家刑名分类，不能等同于统一政治组织随之进入后续年序。《后汉书·桓帝纪》；《后汉书·党锢列传》；《后汉纪》卷九十二提供这一线索，但案件与禁锢存在可靠；被捕人数、党人边界和罪名有争议。
\par

\noindent 在永康元年（168年），\eventanchor{EH-0168-LINGDI-ACCESSION}{桓帝去世，十二岁的刘宏即位，是为灵帝，窦太后、窦武与陈蕃辅政}把东汉置于桓帝无子且宦官势力已深植禁中，朝廷选择河间王刘开一支的刘宏维持刘氏皇统与幼帝即位使外戚、士人大臣与宦官围绕辅政权和禁军控制展开直接竞争之间。可据《后汉书·桓帝纪》；《后汉书·灵帝纪》；《后汉书·窦武传》追踪其大体经过；继位年月可靠；窦武、陈蕃的职权边界有争议。
\par

\noindent 东汉在建宁元年（168年）的\eventanchor{EH-0168-DOU-WU-COUP}{宦官曹节、王甫等发动政变，杀窦武、陈蕃，窦太后被幽禁}，并非孤立一事：窦武、陈蕃谋诛宦官但未能稳控宫门、诏令和北军，宦官利用皇帝近侍地位反击，且外戚—士人联合失败，宦官取得压倒性禁中优势，灵帝朝政治空间急剧收缩。材料见《后汉书·灵帝纪》；《后汉书·窦武传》；《后汉书·陈蕃传》；政变事实可靠；先发制人责任与战斗细节有争议。
\par

\noindent \eventanchor{EH-0169-SECOND-PARTISAN-PROHIBITION}{宦官集团扩大逮捕、禁锢党人，第二次党锢形成}发生于建宁二年（169年），行动者为东汉。它回应窦武政变后宦官将士人批评视为潜在复辟威胁，以结党罪名清除官学和地方声望网络，并把大量士人失去仕进资格，中央与地方精英沟通受损，后期危机中清议人物将以不同方式重返政治留给随后的人群和官署。《后汉书·灵帝纪》；《后汉书·党锢列传》；《资治通鉴》卷五十六可作核验，惟禁锢发生可靠；牵连名单、地域规模和持续期限有争议。
\par

\noindent 从宦官和皇帝试图绕开太学清议与经师网络，培植较少依赖旧式荐举的近侍文化资源出发，东汉于光和元年（178年）出现\eventanchor{EH-0178-HONGDU-GATE}{灵帝置鸿都门学，招集辞赋、书画与尺牍之士入学}；文学与书画技能获得官方位置，也加剧士人对卖官、近幸和选官标准的批评。这里采用《后汉书·灵帝纪》；《后汉书·蔡邕传》；《后汉纪》卷一百零四的记载范围，不越过设置事实可靠；学员身份、课程与政治功能有争议。
\par

\noindent 建安十三年（208年），曹操集团与荆州刘表集团的\eventanchor{EH-0208-JINGZHOU-TRANSFER}{刘表去世后，刘琮降曹，曹操取得荆州北部并南下追击刘备}使刘表长期维持荆州相对独立，病逝后继承集团分裂，曹操以北方优势迅速压迫襄阳转为可见的行动，曹操获得部分荆州军政资源并逼近长江，孙权与刘备面临结盟或各自被吞并的选择。其依据为《三国志·武帝纪》；《三国志·先主传》；《后汉书·刘表传》，刘表卒与刘琮降曹可靠；降附动机、蔡氏作用和荆州实际交接范围有争议。
\par

\noindent \eventanchor{EH-0208-RED-CLIFFS}{孙刘联军在赤壁一带击败曹操水军，曹操撤回北方}这一建安十三年冬（208年）的节点，关联曹操集团与孙权刘备联军、曹操快速南下后水军、疫病与江淮—荆州水域经验不足，孙刘联盟则须阻止其控制长江和曹操未能一举统一南方，孙权和刘备各自取得扩张空间；演义化细节不可替代纪传材料。《三国志·武帝纪》；《三国志·周瑜传》；《三国志·先主传》；《资治通鉴》卷六十五记录了它的基本轮廓；战役结果可靠；战场位置、火攻执行者、疫病影响和兵数有争议。
\par

\noindent 孙权、刘备与曹操集团在建安十四年（209年）处理赤壁后曹操退守北方，孙刘需在共同抗曹与各自获得荆州人口、港口和粮赋之间协商时，出现\eventanchor{EH-0209-JINGZHOU-DIVISION}{孙刘在荆州南部及江陵周边重新分配城郡与防线，三方沿长江形成对峙}；荆州成为孙刘联盟的利益冲突核心，三方边界并非固定，后续战争围绕其交通和郡县展开。可核材料为《三国志·周瑜传》；《三国志·先主传》；《三国志·鲁肃传》，且控制格局大体可靠；“借荆州”条款、具体郡界和交接时间有争议。
\par

\noindent 建安十六年（211年）的\eventanchor{EH-0211-TONG-PASS}{曹操在潼关击败马超、韩遂等关中联军，关中军事集团分裂}，让曹操集团与关中诸将的曹操西进威胁关中诸将的地盘与补给，马超、韩遂等暂时联合以守关中门户在具体事务中显形；关中联军瓦解，曹操取得西向战略主动，但凉州与汉中仍非完全受控区域。《三国志·武帝纪》；《三国志·马超传》；《后汉书·董卓传》所能支持的范围仍受战役结果可靠；离间计、渡河路线和联军构成有争议限制。
\par

\noindent 对刘备集团与益州刘璋政权而言，\eventanchor{EH-0212-LIU-BEI-YIZHOU}{刘备由葭萌进军益州，与刘璋决裂，益州内战开始}发生在建安十七年（212年）：刘璋为防张鲁邀刘备入蜀，却无法协调益州旧有官僚、豪强和外来刘备军的利益，随后刘备获得争夺巴蜀的立足点，益州人口、粮赋与山地关隘成为三方格局的关键资源。参看《三国志·先主传》；《三国志·刘璋传》；《华阳国志·刘先主志》；进军与决裂事实可靠；刘备受邀初衷、庞统建议和开战责任有争议。
\par

\noindent 刘备军控制益州要地并获部分地方官民支持，刘璋难以组织持续反攻或外援之下，刘备集团与益州刘璋政权于建安十九年（214年）经历\eventanchor{EH-0214-CHENGDU-SURRENDER}{刘璋降刘备，刘备据成都并接收益州行政与军需资源}，并面对刘备获得稳固的粮赋与兵员基地，可与曹操、孙权并列竞争；旧益州集团的整合仍须通过任官和联盟完成。此处依《三国志·先主传》；《三国志·刘璋传》；《华阳国志·刘先主志》，降附年份可靠；围城谈判、地方士族选择和接收程度有争议。
\par

\noindent 建安二十年（215年），曹操集团与张鲁政权面对张鲁政权控制汉中多年，曹操在平定关中后需消除其对关中南缘与益州北门的威胁；\eventanchor{EH-0215-CAO-CAO-HANZHONG}{曹操攻取汉中，张鲁降，曹操控制连接关中与巴蜀的北部通道}，其后曹操获得汉中并迁徙部分居民，刘备感到益州北部受压，汉中成为双方下一轮战争焦点。记事依据《三国志·武帝纪》；《三国志·张鲁传》；《后汉书·张鲁传》；攻取与降附事实可靠；张鲁宗教组织、人口迁徙和曹操留守部署有争议。
\par

\noindent 建安二十年（215年）的孙权与刘备集团发生\eventanchor{EH-0215-SUN-QUAN-JINGZHOU}{孙权夺取长沙、零陵、桂阳等荆南郡，孙刘以湘水附近重新划界}。刘备据益州后仍控制部分荆州，孙权认为赤壁后分配失衡并需保障长江上游安全构成其条件，孙刘联盟虽暂时修复以共同对曹，却将荆州主权冲突制度化，关羽日后北进缺少稳固后方随之进入后续年序。《三国志·先主传》；《三国志·孙权传》；《三国志·鲁肃传》提供这一线索，但夺郡与和议大体可靠；各郡先后、湘水划界位置和实际执行有争议。
\par

\noindent 在建安二十二年（217年），\eventanchor{EH-0217-JIANYAN-PLAGUE}{中原发生大疫，曹丕、王粲等文士和军民遭受损失，战争动员承受额外压力}把曹操控制区置于连年征战、人口流动、营垒与灾害使疾病传播风险上升，官府救济与医疗记录有限与文士群体的死亡留下可见文本记录，军政与地方社会的损失难以量化，不能以文学哀悼推算全国死亡率之间。可据《三国志·武帝纪》；曹植《说疫气》；《文选》所收建安文人作品追踪其大体经过；疫病流行可靠；死亡规模、病原和区域范围无法确定。
\par

\noindent 刘备集团与曹操集团在建安二十三年（218年）的\eventanchor{EH-0218-HANZHONG-CAMPAIGN}{刘备在汉中持续进攻夏侯渊、张郃等部，争夺定军山和阳平关方向的通道}，并非孤立一事：曹操占汉中后威胁益州北部，刘备需保住巴蜀门户并取得可防御的山地屏障，且汉中战争消耗双方兵力和粮运，为次年刘备取得汉中及关羽北伐创造相互牵连的战场条件。材料见《三国志·先主传》；《三国志·夏侯渊传》；《三国志·张郃传》；战役阶段可靠；行军路线、各将兵额和地方协力程度有争议。
\par

\noindent \eventanchor{EH-0219-DINGJUN-HANZHONG}{刘备军在定军山击杀夏侯渊，曹操亲至后仍撤出汉中，刘备据汉中称汉中王}发生于建安二十四年（219年），行动者为刘备集团与曹操集团。它回应刘备长期依托益州供给围攻汉中，夏侯渊前出部署使曹军失去关键统帅，并把刘备获得汉中防线与“汉中王”名号，三方竞争从州郡割据进一步转向准王朝建构留给随后的人群和官署。《三国志·先主传》；《三国志·夏侯渊传》；《三国志·法正传》可作核验，惟核心结果可靠；黄忠战功、战场细节和称王礼仪有争议。
\par

\noindent 从关羽北进利用汉水洪灾压迫曹军，却使荆州后方防御薄弱；孙权长期不满荆州分配而转向突袭出发，关羽集团、曹操集团与孙权集团于建安二十四年（219年）出现\eventanchor{EH-0219-FANCHENG-GUAN-YU}{关羽北攻樊城并水淹于禁军，后孙权袭取荆州，关羽兵败被杀}；关羽集团覆灭，刘备失去荆州与东部通道，孙刘同盟实质破裂，曹魏、蜀汉、孙吴三方边界趋于定型。这里采用《三国志·关羽传》；《三国志·于禁传》；《三国志·吕蒙传》；《资治通鉴》卷六十八的记载范围，不越过战役序列可靠；水淹原因、吕蒙白衣渡江细节和关羽被杀地点有争议。
\par

\noindent 建武中元二年（57年），东汉的\eventanchor{EH-LEDGER-013-1}{“光武帝去世，太子刘庄即位，是为明帝”的前期动员与资源准备}使当时的权力、资源与信息约束推动相关过程展开转为可见的行动，这一过程为“光武帝去世，太子刘庄即位，是为明帝”提供可观察的条件。其依据为《后汉书·光武帝纪下》；《后汉书·明帝纪》；《后汉纪》卷二十八，核心事项已有既有账本来源；本分解节点的参与者、执行范围和时间先后仍须逐案核验。
\par

\noindent \eventanchor{EH-LEDGER-013-2}{“光武帝去世，太子刘庄即位，是为明帝”的命令、联盟或地方响应}这一建武中元二年（57年）的节点，关联东汉、当时的权力、资源与信息约束推动相关过程展开和这一过程为“光武帝去世，太子刘庄即位，是为明帝”提供可观察的条件。《后汉书·光武帝纪下》；《后汉书·明帝纪》；《后汉纪》卷二十八记录了它的基本轮廓；核心事项已有既有账本来源；本分解节点的参与者、执行范围和时间先后仍须逐案核验。
\par

\noindent 东汉在建武中元二年（57年）处理“光武帝去世，太子刘庄即位，是为明帝”发生的政治与区域条件共同作用时，出现\eventanchor{EH-LEDGER-013-3}{光武帝去世，太子刘庄即位，是为明帝}；该节点改变后续资源、联盟或制度处置的条件。可核材料为《后汉书·光武帝纪下》；《后汉书·明帝纪》；《后汉纪》卷二十八，且核心事项已有既有账本来源；本分解节点的参与者、执行范围和时间先后仍须逐案核验。
\par

\noindent 建武中元二年（57年）的\eventanchor{EH-LEDGER-013-4}{“光武帝去世，太子刘庄即位，是为明帝”的执行中的空间与人员处置}，让东汉的“光武帝去世，太子刘庄即位，是为明帝”发生的政治与区域条件共同作用在具体事务中显形；该节点改变后续资源、联盟或制度处置的条件。《后汉书·光武帝纪下》；《后汉书·明帝纪》；《后汉纪》卷二十八所能支持的范围仍受核心事项已有既有账本来源；本分解节点的参与者、执行范围和时间先后仍须逐案核验限制。
\par

\noindent 对东汉而言，\eventanchor{EH-LEDGER-013-5}{“光武帝去世，太子刘庄即位，是为明帝”的后续制度或军政调整}发生在建武中元二年（57年）：核心节点发生后仍须处理其遗留问题，随后后续影响须在不同地区和材料类型中分别核验。参看《后汉书·光武帝纪下》；《后汉书·明帝纪》；《后汉纪》卷二十八；核心事项已有既有账本来源；本分解节点的参与者、执行范围和时间先后仍须逐案核验。
\par

\noindent 核心节点发生后仍须处理其遗留问题之下，东汉于建武中元二年（57年）经历\eventanchor{EH-LEDGER-013-6}{“光武帝去世，太子刘庄即位，是为明帝”的异源材料与影响范围的复核}，并面对后续影响须在不同地区和材料类型中分别核验。此处依《后汉书·光武帝纪下》；《后汉书·明帝纪》；《后汉纪》卷二十八，核心事项已有既有账本来源；本分解节点的参与者、执行范围和时间先后仍须逐案核验。
\par

\noindent 永平八年（65年），东汉面对当时的权力、资源与信息约束推动相关过程展开；\eventanchor{EH-LEDGER-014-1}{“楚王刘英因被告谋反获减死并被徙，其诏书涉及黄老与佛教供养的早期记录”的前期动员与资源准备}，其后这一过程为“楚王刘英因被告谋反获减死并被徙，其诏书涉及黄老与佛教供养的早期记录”提供可观察的条件。记事依据《后汉书·楚王英传》；《后汉书·明帝纪》；《牟子理惑论》相关传本；核心事项已有既有账本来源；本分解节点的参与者、执行范围和时间先后仍须逐案核验。
\par

\noindent 永平八年（65年）的东汉发生\eventanchor{EH-LEDGER-014-2}{“楚王刘英因被告谋反获减死并被徙，其诏书涉及黄老与佛教供养的早期记录”的命令、联盟或地方响应}。当时的权力、资源与信息约束推动相关过程展开构成其条件，这一过程为“楚王刘英因被告谋反获减死并被徙，其诏书涉及黄老与佛教供养的早期记录”提供可观察的条件随之进入后续年序。《后汉书·楚王英传》；《后汉书·明帝纪》；《牟子理惑论》相关传本提供这一线索，但核心事项已有既有账本来源；本分解节点的参与者、执行范围和时间先后仍须逐案核验。
\par

\noindent 在永平八年（65年），\eventanchor{EH-LEDGER-014-3}{楚王刘英因被告谋反获减死并被徙，其诏书涉及黄老与佛教供养的早期记录}把东汉置于“楚王刘英因被告谋反获减死并被徙，其诏书涉及黄老与佛教供养的早期记录”发生的政治与区域条件共同作用与该节点改变后续资源、联盟或制度处置的条件之间。可据《后汉书·楚王英传》；《后汉书·明帝纪》；《牟子理惑论》相关传本追踪其大体经过；核心事项已有既有账本来源；本分解节点的参与者、执行范围和时间先后仍须逐案核验。
\par

\noindent 东汉在永平八年（65年）的\eventanchor{EH-LEDGER-014-4}{“楚王刘英因被告谋反获减死并被徙，其诏书涉及黄老与佛教供养的早期记录”的执行中的空间与人员处置}，并非孤立一事：“楚王刘英因被告谋反获减死并被徙，其诏书涉及黄老与佛教供养的早期记录”发生的政治与区域条件共同作用，且该节点改变后续资源、联盟或制度处置的条件。材料见《后汉书·楚王英传》；《后汉书·明帝纪》；《牟子理惑论》相关传本；核心事项已有既有账本来源；本分解节点的参与者、执行范围和时间先后仍须逐案核验。
\par

\noindent \eventanchor{EH-LEDGER-014-5}{“楚王刘英因被告谋反获减死并被徙，其诏书涉及黄老与佛教供养的早期记录”的后续制度或军政调整}发生于永平八年（65年），行动者为东汉。它回应核心节点发生后仍须处理其遗留问题，并把后续影响须在不同地区和材料类型中分别核验留给随后的人群和官署。《后汉书·楚王英传》；《后汉书·明帝纪》；《牟子理惑论》相关传本可作核验，惟核心事项已有既有账本来源；本分解节点的参与者、执行范围和时间先后仍须逐案核验。
\par

\noindent 从核心节点发生后仍须处理其遗留问题出发，东汉于永平八年（65年）出现\eventanchor{EH-LEDGER-014-6}{“楚王刘英因被告谋反获减死并被徙，其诏书涉及黄老与佛教供养的早期记录”的异源材料与影响范围的复核}；后续影响须在不同地区和材料类型中分别核验。这里采用《后汉书·楚王英传》；《后汉书·明帝纪》；《牟子理惑论》相关传本的记载范围，不越过核心事项已有既有账本来源；本分解节点的参与者、执行范围和时间先后仍须逐案核验。
\par

\noindent 永平十八年（75年），东汉的\eventanchor{EH-LEDGER-017-1}{“明帝去世，太子刘炟即位，是为章帝”的前期动员与资源准备}使当时的权力、资源与信息约束推动相关过程展开转为可见的行动，这一过程为“明帝去世，太子刘炟即位，是为章帝”提供可观察的条件。其依据为《后汉书·明帝纪》；《后汉书·章帝纪》；《后汉纪》卷三十六，核心事项已有既有账本来源；本分解节点的参与者、执行范围和时间先后仍须逐案核验。
\par

\noindent \eventanchor{EH-LEDGER-017-2}{“明帝去世，太子刘炟即位，是为章帝”的命令、联盟或地方响应}这一永平十八年（75年）的节点，关联东汉、当时的权力、资源与信息约束推动相关过程展开和这一过程为“明帝去世，太子刘炟即位，是为章帝”提供可观察的条件。《后汉书·明帝纪》；《后汉书·章帝纪》；《后汉纪》卷三十六记录了它的基本轮廓；核心事项已有既有账本来源；本分解节点的参与者、执行范围和时间先后仍须逐案核验。
\par

\noindent 东汉在永平十八年（75年）处理“明帝去世，太子刘炟即位，是为章帝”发生的政治与区域条件共同作用时，出现\eventanchor{EH-LEDGER-017-3}{明帝去世，太子刘炟即位，是为章帝}；该节点改变后续资源、联盟或制度处置的条件。可核材料为《后汉书·明帝纪》；《后汉书·章帝纪》；《后汉纪》卷三十六，且核心事项已有既有账本来源；本分解节点的参与者、执行范围和时间先后仍须逐案核验。
\par

\noindent 永平十八年（75年）的\eventanchor{EH-LEDGER-017-4}{“明帝去世，太子刘炟即位，是为章帝”的执行中的空间与人员处置}，让东汉的“明帝去世，太子刘炟即位，是为章帝”发生的政治与区域条件共同作用在具体事务中显形；该节点改变后续资源、联盟或制度处置的条件。《后汉书·明帝纪》；《后汉书·章帝纪》；《后汉纪》卷三十六所能支持的范围仍受核心事项已有既有账本来源；本分解节点的参与者、执行范围和时间先后仍须逐案核验限制。
\par

\noindent 对东汉而言，\eventanchor{EH-LEDGER-017-5}{“明帝去世，太子刘炟即位，是为章帝”的后续制度或军政调整}发生在永平十八年（75年）：核心节点发生后仍须处理其遗留问题，随后后续影响须在不同地区和材料类型中分别核验。参看《后汉书·明帝纪》；《后汉书·章帝纪》；《后汉纪》卷三十六；核心事项已有既有账本来源；本分解节点的参与者、执行范围和时间先后仍须逐案核验。
\par

\noindent 核心节点发生后仍须处理其遗留问题之下，东汉于永平十八年（75年）经历\eventanchor{EH-LEDGER-017-6}{“明帝去世，太子刘炟即位，是为章帝”的异源材料与影响范围的复核}，并面对后续影响须在不同地区和材料类型中分别核验。此处依《后汉书·明帝纪》；《后汉书·章帝纪》；《后汉纪》卷三十六，核心事项已有既有账本来源；本分解节点的参与者、执行范围和时间先后仍须逐案核验。
\par

\noindent 建初四年（79年），东汉面对当时的权力、资源与信息约束推动相关过程展开；\eventanchor{EH-LEDGER-019-1}{“章帝召集儒者在白虎观讨论五经异同，形成《白虎通》所保存的经义框架”的前期动员与资源准备}，其后这一过程为“章帝召集儒者在白虎观讨论五经异同，形成《白虎通》所保存的经义框架”提供可观察的条件。记事依据《后汉书·章帝纪》；《白虎通义》；《后汉纪》卷四十；核心事项已有既有账本来源；本分解节点的参与者、执行范围和时间先后仍须逐案核验。
\par

\noindent 建初四年（79年）的东汉发生\eventanchor{EH-LEDGER-019-2}{“章帝召集儒者在白虎观讨论五经异同，形成《白虎通》所保存的经义框架”的命令、联盟或地方响应}。当时的权力、资源与信息约束推动相关过程展开构成其条件，这一过程为“章帝召集儒者在白虎观讨论五经异同，形成《白虎通》所保存的经义框架”提供可观察的条件随之进入后续年序。《后汉书·章帝纪》；《白虎通义》；《后汉纪》卷四十提供这一线索，但核心事项已有既有账本来源；本分解节点的参与者、执行范围和时间先后仍须逐案核验。
\par

\noindent 在建初四年（79年），\eventanchor{EH-LEDGER-019-3}{章帝召集儒者在白虎观讨论五经异同，形成《白虎通》所保存的经义框架}把东汉置于“章帝召集儒者在白虎观讨论五经异同，形成《白虎通》所保存的经义框架”发生的政治与区域条件共同作用与该节点改变后续资源、联盟或制度处置的条件之间。可据《后汉书·章帝纪》；《白虎通义》；《后汉纪》卷四十追踪其大体经过；核心事项已有既有账本来源；本分解节点的参与者、执行范围和时间先后仍须逐案核验。
\par

\noindent 东汉在建初四年（79年）的\eventanchor{EH-LEDGER-019-4}{“章帝召集儒者在白虎观讨论五经异同，形成《白虎通》所保存的经义框架”的执行中的空间与人员处置}，并非孤立一事：“章帝召集儒者在白虎观讨论五经异同，形成《白虎通》所保存的经义框架”发生的政治与区域条件共同作用，且该节点改变后续资源、联盟或制度处置的条件。材料见《后汉书·章帝纪》；《白虎通义》；《后汉纪》卷四十；核心事项已有既有账本来源；本分解节点的参与者、执行范围和时间先后仍须逐案核验。
\par

\noindent \eventanchor{EH-LEDGER-019-5}{“章帝召集儒者在白虎观讨论五经异同，形成《白虎通》所保存的经义框架”的后续制度或军政调整}发生于建初四年（79年），行动者为东汉。它回应核心节点发生后仍须处理其遗留问题，并把后续影响须在不同地区和材料类型中分别核验留给随后的人群和官署。《后汉书·章帝纪》；《白虎通义》；《后汉纪》卷四十可作核验，惟核心事项已有既有账本来源；本分解节点的参与者、执行范围和时间先后仍须逐案核验。
\par

\noindent 从核心节点发生后仍须处理其遗留问题出发，东汉于建初四年（79年）出现\eventanchor{EH-LEDGER-019-6}{“章帝召集儒者在白虎观讨论五经异同，形成《白虎通》所保存的经义框架”的异源材料与影响范围的复核}；后续影响须在不同地区和材料类型中分别核验。这里采用《后汉书·章帝纪》；《白虎通义》；《后汉纪》卷四十的记载范围，不越过核心事项已有既有账本来源；本分解节点的参与者、执行范围和时间先后仍须逐案核验。
\par

\noindent 章和二年（88年），东汉的\eventanchor{EH-LEDGER-020-1}{“章帝去世，皇子刘肇即位，是为和帝，窦太后临朝”的前期动员与资源准备}使当时的权力、资源与信息约束推动相关过程展开转为可见的行动，这一过程为“章帝去世，皇子刘肇即位，是为和帝，窦太后临朝”提供可观察的条件。其依据为《后汉书·章帝纪》；《后汉书·和帝纪》；《后汉纪》卷四十九，核心事项已有既有账本来源；本分解节点的参与者、执行范围和时间先后仍须逐案核验。
\par

\noindent \eventanchor{EH-LEDGER-020-2}{“章帝去世，皇子刘肇即位，是为和帝，窦太后临朝”的命令、联盟或地方响应}这一章和二年（88年）的节点，关联东汉、当时的权力、资源与信息约束推动相关过程展开和这一过程为“章帝去世，皇子刘肇即位，是为和帝，窦太后临朝”提供可观察的条件。《后汉书·章帝纪》；《后汉书·和帝纪》；《后汉纪》卷四十九记录了它的基本轮廓；核心事项已有既有账本来源；本分解节点的参与者、执行范围和时间先后仍须逐案核验。
\par

\noindent 东汉在章和二年（88年）处理“章帝去世，皇子刘肇即位，是为和帝，窦太后临朝”发生的政治与区域条件共同作用时，出现\eventanchor{EH-LEDGER-020-3}{章帝去世，皇子刘肇即位，是为和帝，窦太后临朝}；该节点改变后续资源、联盟或制度处置的条件。可核材料为《后汉书·章帝纪》；《后汉书·和帝纪》；《后汉纪》卷四十九，且核心事项已有既有账本来源；本分解节点的参与者、执行范围和时间先后仍须逐案核验。
\par

\noindent 章和二年（88年）的\eventanchor{EH-LEDGER-020-4}{“章帝去世，皇子刘肇即位，是为和帝，窦太后临朝”的执行中的空间与人员处置}，让东汉的“章帝去世，皇子刘肇即位，是为和帝，窦太后临朝”发生的政治与区域条件共同作用在具体事务中显形；该节点改变后续资源、联盟或制度处置的条件。《后汉书·章帝纪》；《后汉书·和帝纪》；《后汉纪》卷四十九所能支持的范围仍受核心事项已有既有账本来源；本分解节点的参与者、执行范围和时间先后仍须逐案核验限制。
\par

\noindent 对东汉而言，\eventanchor{EH-LEDGER-020-5}{“章帝去世，皇子刘肇即位，是为和帝，窦太后临朝”的后续制度或军政调整}发生在章和二年（88年）：核心节点发生后仍须处理其遗留问题，随后后续影响须在不同地区和材料类型中分别核验。参看《后汉书·章帝纪》；《后汉书·和帝纪》；《后汉纪》卷四十九；核心事项已有既有账本来源；本分解节点的参与者、执行范围和时间先后仍须逐案核验。
\par

\noindent 核心节点发生后仍须处理其遗留问题之下，东汉于章和二年（88年）经历\eventanchor{EH-LEDGER-020-6}{“章帝去世，皇子刘肇即位，是为和帝，窦太后临朝”的异源材料与影响范围的复核}，并面对后续影响须在不同地区和材料类型中分别核验。此处依《后汉书·章帝纪》；《后汉书·和帝纪》；《后汉纪》卷四十九，核心事项已有既有账本来源；本分解节点的参与者、执行范围和时间先后仍须逐案核验。
\par

\noindent 永元四年（92年），东汉面对当时的权力、资源与信息约束推动相关过程展开；\eventanchor{EH-LEDGER-023-1}{“和帝联同宦官郑众诛除窦宪集团，结束窦太后家族的军政主导”的前期动员与资源准备}，其后这一过程为“和帝联同宦官郑众诛除窦宪集团，结束窦太后家族的军政主导”提供可观察的条件。记事依据《后汉书·和帝纪》；《后汉书·窦宪传》；《后汉书·宦者列传》；核心事项已有既有账本来源；本分解节点的参与者、执行范围和时间先后仍须逐案核验。
\par

\noindent 永元四年（92年）的东汉发生\eventanchor{EH-LEDGER-023-2}{“和帝联同宦官郑众诛除窦宪集团，结束窦太后家族的军政主导”的命令、联盟或地方响应}。当时的权力、资源与信息约束推动相关过程展开构成其条件，这一过程为“和帝联同宦官郑众诛除窦宪集团，结束窦太后家族的军政主导”提供可观察的条件随之进入后续年序。《后汉书·和帝纪》；《后汉书·窦宪传》；《后汉书·宦者列传》提供这一线索，但核心事项已有既有账本来源；本分解节点的参与者、执行范围和时间先后仍须逐案核验。
\par

\noindent 在永元四年（92年），\eventanchor{EH-LEDGER-023-3}{和帝联同宦官郑众诛除窦宪集团，结束窦太后家族的军政主导}把东汉置于“和帝联同宦官郑众诛除窦宪集团，结束窦太后家族的军政主导”发生的政治与区域条件共同作用与该节点改变后续资源、联盟或制度处置的条件之间。可据《后汉书·和帝纪》；《后汉书·窦宪传》；《后汉书·宦者列传》追踪其大体经过；核心事项已有既有账本来源；本分解节点的参与者、执行范围和时间先后仍须逐案核验。
\par

\noindent 东汉在永元四年（92年）的\eventanchor{EH-LEDGER-023-4}{“和帝联同宦官郑众诛除窦宪集团，结束窦太后家族的军政主导”的执行中的空间与人员处置}，并非孤立一事：“和帝联同宦官郑众诛除窦宪集团，结束窦太后家族的军政主导”发生的政治与区域条件共同作用，且该节点改变后续资源、联盟或制度处置的条件。材料见《后汉书·和帝纪》；《后汉书·窦宪传》；《后汉书·宦者列传》；核心事项已有既有账本来源；本分解节点的参与者、执行范围和时间先后仍须逐案核验。
\par

\noindent \eventanchor{EH-LEDGER-023-5}{“和帝联同宦官郑众诛除窦宪集团，结束窦太后家族的军政主导”的后续制度或军政调整}发生于永元四年（92年），行动者为东汉。它回应核心节点发生后仍须处理其遗留问题，并把后续影响须在不同地区和材料类型中分别核验留给随后的人群和官署。《后汉书·和帝纪》；《后汉书·窦宪传》；《后汉书·宦者列传》可作核验，惟核心事项已有既有账本来源；本分解节点的参与者、执行范围和时间先后仍须逐案核验。
\par

\noindent 从核心节点发生后仍须处理其遗留问题出发，东汉于永元四年（92年）出现\eventanchor{EH-LEDGER-023-6}{“和帝联同宦官郑众诛除窦宪集团，结束窦太后家族的军政主导”的异源材料与影响范围的复核}；后续影响须在不同地区和材料类型中分别核验。这里采用《后汉书·和帝纪》；《后汉书·窦宪传》；《后汉书·宦者列传》的记载范围，不越过核心事项已有既有账本来源；本分解节点的参与者、执行范围和时间先后仍须逐案核验。
\par

\noindent 元兴元年（105年），东汉的\eventanchor{EH-LEDGER-026-1}{“和帝去世，襁褓中的刘隆即位，是为殇帝，邓太后临朝”的前期动员与资源准备}使当时的权力、资源与信息约束推动相关过程展开转为可见的行动，这一过程为“和帝去世，襁褓中的刘隆即位，是为殇帝，邓太后临朝”提供可观察的条件。其依据为《后汉书·和帝纪》；《后汉书·殇帝纪》；《后汉书·邓皇后纪》，核心事项已有既有账本来源；本分解节点的参与者、执行范围和时间先后仍须逐案核验。
\par

\noindent \eventanchor{EH-LEDGER-026-2}{“和帝去世，襁褓中的刘隆即位，是为殇帝，邓太后临朝”的命令、联盟或地方响应}这一元兴元年（105年）的节点，关联东汉、当时的权力、资源与信息约束推动相关过程展开和这一过程为“和帝去世，襁褓中的刘隆即位，是为殇帝，邓太后临朝”提供可观察的条件。《后汉书·和帝纪》；《后汉书·殇帝纪》；《后汉书·邓皇后纪》记录了它的基本轮廓；核心事项已有既有账本来源；本分解节点的参与者、执行范围和时间先后仍须逐案核验。
\par

\noindent 东汉在元兴元年（105年）处理“和帝去世，襁褓中的刘隆即位，是为殇帝，邓太后临朝”发生的政治与区域条件共同作用时，出现\eventanchor{EH-LEDGER-026-3}{和帝去世，襁褓中的刘隆即位，是为殇帝，邓太后临朝}；该节点改变后续资源、联盟或制度处置的条件。可核材料为《后汉书·和帝纪》；《后汉书·殇帝纪》；《后汉书·邓皇后纪》，且核心事项已有既有账本来源；本分解节点的参与者、执行范围和时间先后仍须逐案核验。
\par

\noindent 元兴元年（105年）的\eventanchor{EH-LEDGER-026-4}{“和帝去世，襁褓中的刘隆即位，是为殇帝，邓太后临朝”的执行中的空间与人员处置}，让东汉的“和帝去世，襁褓中的刘隆即位，是为殇帝，邓太后临朝”发生的政治与区域条件共同作用在具体事务中显形；该节点改变后续资源、联盟或制度处置的条件。《后汉书·和帝纪》；《后汉书·殇帝纪》；《后汉书·邓皇后纪》所能支持的范围仍受核心事项已有既有账本来源；本分解节点的参与者、执行范围和时间先后仍须逐案核验限制。
\par

\noindent 对东汉而言，\eventanchor{EH-LEDGER-026-5}{“和帝去世，襁褓中的刘隆即位，是为殇帝，邓太后临朝”的后续制度或军政调整}发生在元兴元年（105年）：核心节点发生后仍须处理其遗留问题，随后后续影响须在不同地区和材料类型中分别核验。参看《后汉书·和帝纪》；《后汉书·殇帝纪》；《后汉书·邓皇后纪》；核心事项已有既有账本来源；本分解节点的参与者、执行范围和时间先后仍须逐案核验。
\par

\noindent 核心节点发生后仍须处理其遗留问题之下，东汉于元兴元年（105年）经历\eventanchor{EH-LEDGER-026-6}{“和帝去世，襁褓中的刘隆即位，是为殇帝，邓太后临朝”的异源材料与影响范围的复核}，并面对后续影响须在不同地区和材料类型中分别核验。此处依《后汉书·和帝纪》；《后汉书·殇帝纪》；《后汉书·邓皇后纪》，核心事项已有既有账本来源；本分解节点的参与者、执行范围和时间先后仍须逐案核验。
\par

\noindent 延平元年（106年），东汉面对当时的权力、资源与信息约束推动相关过程展开；\eventanchor{EH-LEDGER-027-1}{“殇帝去世，清河王刘祜即位，是为安帝，邓太后继续临朝”的前期动员与资源准备}，其后这一过程为“殇帝去世，清河王刘祜即位，是为安帝，邓太后继续临朝”提供可观察的条件。记事依据《后汉书·殇帝纪》；《后汉书·安帝纪》；《后汉书·邓皇后纪》；核心事项已有既有账本来源；本分解节点的参与者、执行范围和时间先后仍须逐案核验。
\par

\noindent 延平元年（106年）的东汉发生\eventanchor{EH-LEDGER-027-2}{“殇帝去世，清河王刘祜即位，是为安帝，邓太后继续临朝”的命令、联盟或地方响应}。当时的权力、资源与信息约束推动相关过程展开构成其条件，这一过程为“殇帝去世，清河王刘祜即位，是为安帝，邓太后继续临朝”提供可观察的条件随之进入后续年序。《后汉书·殇帝纪》；《后汉书·安帝纪》；《后汉书·邓皇后纪》提供这一线索，但核心事项已有既有账本来源；本分解节点的参与者、执行范围和时间先后仍须逐案核验。
\par

\noindent 在延平元年（106年），\eventanchor{EH-LEDGER-027-3}{殇帝去世，清河王刘祜即位，是为安帝，邓太后继续临朝}把东汉置于“殇帝去世，清河王刘祜即位，是为安帝，邓太后继续临朝”发生的政治与区域条件共同作用与该节点改变后续资源、联盟或制度处置的条件之间。可据《后汉书·殇帝纪》；《后汉书·安帝纪》；《后汉书·邓皇后纪》追踪其大体经过；核心事项已有既有账本来源；本分解节点的参与者、执行范围和时间先后仍须逐案核验。
\par

\noindent 东汉在延平元年（106年）的\eventanchor{EH-LEDGER-027-4}{“殇帝去世，清河王刘祜即位，是为安帝，邓太后继续临朝”的执行中的空间与人员处置}，并非孤立一事：“殇帝去世，清河王刘祜即位，是为安帝，邓太后继续临朝”发生的政治与区域条件共同作用，且该节点改变后续资源、联盟或制度处置的条件。材料见《后汉书·殇帝纪》；《后汉书·安帝纪》；《后汉书·邓皇后纪》；核心事项已有既有账本来源；本分解节点的参与者、执行范围和时间先后仍须逐案核验。
\par

\noindent \eventanchor{EH-LEDGER-027-5}{“殇帝去世，清河王刘祜即位，是为安帝，邓太后继续临朝”的后续制度或军政调整}发生于延平元年（106年），行动者为东汉。它回应核心节点发生后仍须处理其遗留问题，并把后续影响须在不同地区和材料类型中分别核验留给随后的人群和官署。《后汉书·殇帝纪》；《后汉书·安帝纪》；《后汉书·邓皇后纪》可作核验，惟核心事项已有既有账本来源；本分解节点的参与者、执行范围和时间先后仍须逐案核验。
\par

\noindent 从核心节点发生后仍须处理其遗留问题出发，东汉于延平元年（106年）出现\eventanchor{EH-LEDGER-027-6}{“殇帝去世，清河王刘祜即位，是为安帝，邓太后继续临朝”的异源材料与影响范围的复核}；后续影响须在不同地区和材料类型中分别核验。这里采用《后汉书·殇帝纪》；《后汉书·安帝纪》；《后汉书·邓皇后纪》的记载范围，不越过核心事项已有既有账本来源；本分解节点的参与者、执行范围和时间先后仍须逐案核验。
\par

\noindent 元初三年（116年），东汉的\eventanchor{EH-LEDGER-029-1}{“邓太后去世，安帝开始亲政并贬逐部分邓氏成员”的前期动员与资源准备}使当时的权力、资源与信息约束推动相关过程展开转为可见的行动，这一过程为“邓太后去世，安帝开始亲政并贬逐部分邓氏成员”提供可观察的条件。其依据为《后汉书·安帝纪》；《后汉书·邓皇后纪》；《后汉纪》卷七十七，核心事项已有既有账本来源；本分解节点的参与者、执行范围和时间先后仍须逐案核验。
\par

\noindent \eventanchor{EH-LEDGER-029-2}{“邓太后去世，安帝开始亲政并贬逐部分邓氏成员”的命令、联盟或地方响应}这一元初三年（116年）的节点，关联东汉、当时的权力、资源与信息约束推动相关过程展开和这一过程为“邓太后去世，安帝开始亲政并贬逐部分邓氏成员”提供可观察的条件。《后汉书·安帝纪》；《后汉书·邓皇后纪》；《后汉纪》卷七十七记录了它的基本轮廓；核心事项已有既有账本来源；本分解节点的参与者、执行范围和时间先后仍须逐案核验。
\par

\noindent 东汉在元初三年（116年）处理“邓太后去世，安帝开始亲政并贬逐部分邓氏成员”发生的政治与区域条件共同作用时，出现\eventanchor{EH-LEDGER-029-3}{邓太后去世，安帝开始亲政并贬逐部分邓氏成员}；该节点改变后续资源、联盟或制度处置的条件。可核材料为《后汉书·安帝纪》；《后汉书·邓皇后纪》；《后汉纪》卷七十七，且核心事项已有既有账本来源；本分解节点的参与者、执行范围和时间先后仍须逐案核验。
\par

\noindent 元初三年（116年）的\eventanchor{EH-LEDGER-029-4}{“邓太后去世，安帝开始亲政并贬逐部分邓氏成员”的执行中的空间与人员处置}，让东汉的“邓太后去世，安帝开始亲政并贬逐部分邓氏成员”发生的政治与区域条件共同作用在具体事务中显形；该节点改变后续资源、联盟或制度处置的条件。《后汉书·安帝纪》；《后汉书·邓皇后纪》；《后汉纪》卷七十七所能支持的范围仍受核心事项已有既有账本来源；本分解节点的参与者、执行范围和时间先后仍须逐案核验限制。
\par

\noindent 对东汉而言，\eventanchor{EH-LEDGER-029-5}{“邓太后去世，安帝开始亲政并贬逐部分邓氏成员”的后续制度或军政调整}发生在元初三年（116年）：核心节点发生后仍须处理其遗留问题，随后后续影响须在不同地区和材料类型中分别核验。参看《后汉书·安帝纪》；《后汉书·邓皇后纪》；《后汉纪》卷七十七；核心事项已有既有账本来源；本分解节点的参与者、执行范围和时间先后仍须逐案核验。
\par

\noindent 核心节点发生后仍须处理其遗留问题之下，东汉于元初三年（116年）经历\eventanchor{EH-LEDGER-029-6}{“邓太后去世，安帝开始亲政并贬逐部分邓氏成员”的异源材料与影响范围的复核}，并面对后续影响须在不同地区和材料类型中分别核验。此处依《后汉书·安帝纪》；《后汉书·邓皇后纪》；《后汉纪》卷七十七，核心事项已有既有账本来源；本分解节点的参与者、执行范围和时间先后仍须逐案核验。
\par

\noindent 延光四年（125年），东汉面对当时的权力、资源与信息约束推动相关过程展开；\eventanchor{EH-LEDGER-030-1}{“安帝去世，阎后先立北乡侯，宦官孙程等旋迎立济阴王刘保为顺帝”的前期动员与资源准备}，其后这一过程为“安帝去世，阎后先立北乡侯，宦官孙程等旋迎立济阴王刘保为顺帝”提供可观察的条件。记事依据《后汉书·安帝纪》；《后汉书·顺帝纪》；《后汉书·孙程传》；核心事项已有既有账本来源；本分解节点的参与者、执行范围和时间先后仍须逐案核验。
\par

\noindent 延光四年（125年）的东汉发生\eventanchor{EH-LEDGER-030-2}{“安帝去世，阎后先立北乡侯，宦官孙程等旋迎立济阴王刘保为顺帝”的命令、联盟或地方响应}。当时的权力、资源与信息约束推动相关过程展开构成其条件，这一过程为“安帝去世，阎后先立北乡侯，宦官孙程等旋迎立济阴王刘保为顺帝”提供可观察的条件随之进入后续年序。《后汉书·安帝纪》；《后汉书·顺帝纪》；《后汉书·孙程传》提供这一线索，但核心事项已有既有账本来源；本分解节点的参与者、执行范围和时间先后仍须逐案核验。
\par

\noindent 在延光四年（125年），\eventanchor{EH-LEDGER-030-3}{安帝去世，阎后先立北乡侯，宦官孙程等旋迎立济阴王刘保为顺帝}把东汉置于“安帝去世，阎后先立北乡侯，宦官孙程等旋迎立济阴王刘保为顺帝”发生的政治与区域条件共同作用与该节点改变后续资源、联盟或制度处置的条件之间。可据《后汉书·安帝纪》；《后汉书·顺帝纪》；《后汉书·孙程传》追踪其大体经过；核心事项已有既有账本来源；本分解节点的参与者、执行范围和时间先后仍须逐案核验。
\par

\noindent 东汉在延光四年（125年）的\eventanchor{EH-LEDGER-030-4}{“安帝去世，阎后先立北乡侯，宦官孙程等旋迎立济阴王刘保为顺帝”的执行中的空间与人员处置}，并非孤立一事：“安帝去世，阎后先立北乡侯，宦官孙程等旋迎立济阴王刘保为顺帝”发生的政治与区域条件共同作用，且该节点改变后续资源、联盟或制度处置的条件。材料见《后汉书·安帝纪》；《后汉书·顺帝纪》；《后汉书·孙程传》；核心事项已有既有账本来源；本分解节点的参与者、执行范围和时间先后仍须逐案核验。
\par

\noindent \eventanchor{EH-LEDGER-030-5}{“安帝去世，阎后先立北乡侯，宦官孙程等旋迎立济阴王刘保为顺帝”的后续制度或军政调整}发生于延光四年（125年），行动者为东汉。它回应核心节点发生后仍须处理其遗留问题，并把后续影响须在不同地区和材料类型中分别核验留给随后的人群和官署。《后汉书·安帝纪》；《后汉书·顺帝纪》；《后汉书·孙程传》可作核验，惟核心事项已有既有账本来源；本分解节点的参与者、执行范围和时间先后仍须逐案核验。
\par

\noindent 从核心节点发生后仍须处理其遗留问题出发，东汉于延光四年（125年）出现\eventanchor{EH-LEDGER-030-6}{“安帝去世，阎后先立北乡侯，宦官孙程等旋迎立济阴王刘保为顺帝”的异源材料与影响范围的复核}；后续影响须在不同地区和材料类型中分别核验。这里采用《后汉书·安帝纪》；《后汉书·顺帝纪》；《后汉书·孙程传》的记载范围，不越过核心事项已有既有账本来源；本分解节点的参与者、执行范围和时间先后仍须逐案核验。
\par

\noindent 建康元年（144年），东汉的\eventanchor{EH-LEDGER-031-1}{“顺帝去世，幼帝刘炳即位，是为冲帝，梁太后与梁冀临朝”的前期动员与资源准备}使当时的权力、资源与信息约束推动相关过程展开转为可见的行动，这一过程为“顺帝去世，幼帝刘炳即位，是为冲帝，梁太后与梁冀临朝”提供可观察的条件。其依据为《后汉书·顺帝纪》；《后汉书·冲帝纪》；《后汉书·梁冀传》，核心事项已有既有账本来源；本分解节点的参与者、执行范围和时间先后仍须逐案核验。
\par

\noindent \eventanchor{EH-LEDGER-031-2}{“顺帝去世，幼帝刘炳即位，是为冲帝，梁太后与梁冀临朝”的命令、联盟或地方响应}这一建康元年（144年）的节点，关联东汉、当时的权力、资源与信息约束推动相关过程展开和这一过程为“顺帝去世，幼帝刘炳即位，是为冲帝，梁太后与梁冀临朝”提供可观察的条件。《后汉书·顺帝纪》；《后汉书·冲帝纪》；《后汉书·梁冀传》记录了它的基本轮廓；核心事项已有既有账本来源；本分解节点的参与者、执行范围和时间先后仍须逐案核验。
\par

\noindent 东汉在建康元年（144年）处理“顺帝去世，幼帝刘炳即位，是为冲帝，梁太后与梁冀临朝”发生的政治与区域条件共同作用时，出现\eventanchor{EH-LEDGER-031-3}{顺帝去世，幼帝刘炳即位，是为冲帝，梁太后与梁冀临朝}；该节点改变后续资源、联盟或制度处置的条件。可核材料为《后汉书·顺帝纪》；《后汉书·冲帝纪》；《后汉书·梁冀传》，且核心事项已有既有账本来源；本分解节点的参与者、执行范围和时间先后仍须逐案核验。
\par

\noindent 建康元年（144年）的\eventanchor{EH-LEDGER-031-4}{“顺帝去世，幼帝刘炳即位，是为冲帝，梁太后与梁冀临朝”的执行中的空间与人员处置}，让东汉的“顺帝去世，幼帝刘炳即位，是为冲帝，梁太后与梁冀临朝”发生的政治与区域条件共同作用在具体事务中显形；该节点改变后续资源、联盟或制度处置的条件。《后汉书·顺帝纪》；《后汉书·冲帝纪》；《后汉书·梁冀传》所能支持的范围仍受核心事项已有既有账本来源；本分解节点的参与者、执行范围和时间先后仍须逐案核验限制。
\par

\noindent 对东汉而言，\eventanchor{EH-LEDGER-031-5}{“顺帝去世，幼帝刘炳即位，是为冲帝，梁太后与梁冀临朝”的后续制度或军政调整}发生在建康元年（144年）：核心节点发生后仍须处理其遗留问题，随后后续影响须在不同地区和材料类型中分别核验。参看《后汉书·顺帝纪》；《后汉书·冲帝纪》；《后汉书·梁冀传》；核心事项已有既有账本来源；本分解节点的参与者、执行范围和时间先后仍须逐案核验。
\par

\noindent 核心节点发生后仍须处理其遗留问题之下，东汉于建康元年（144年）经历\eventanchor{EH-LEDGER-031-6}{“顺帝去世，幼帝刘炳即位，是为冲帝，梁太后与梁冀临朝”的异源材料与影响范围的复核}，并面对后续影响须在不同地区和材料类型中分别核验。此处依《后汉书·顺帝纪》；《后汉书·冲帝纪》；《后汉书·梁冀传》，核心事项已有既有账本来源；本分解节点的参与者、执行范围和时间先后仍须逐案核验。
\par

\noindent 永嘉元年（145年），东汉面对当时的权力、资源与信息约束推动相关过程展开；\eventanchor{EH-LEDGER-032-1}{“冲帝夭亡，梁太后立八岁的刘缵为质帝，梁冀继续掌政”的前期动员与资源准备}，其后这一过程为“冲帝夭亡，梁太后立八岁的刘缵为质帝，梁冀继续掌政”提供可观察的条件。记事依据《后汉书·冲帝纪》；《后汉书·质帝纪》；《后汉书·梁冀传》；核心事项已有既有账本来源；本分解节点的参与者、执行范围和时间先后仍须逐案核验。
\par

\noindent 永嘉元年（145年）的东汉发生\eventanchor{EH-LEDGER-032-2}{“冲帝夭亡，梁太后立八岁的刘缵为质帝，梁冀继续掌政”的命令、联盟或地方响应}。当时的权力、资源与信息约束推动相关过程展开构成其条件，这一过程为“冲帝夭亡，梁太后立八岁的刘缵为质帝，梁冀继续掌政”提供可观察的条件随之进入后续年序。《后汉书·冲帝纪》；《后汉书·质帝纪》；《后汉书·梁冀传》提供这一线索，但核心事项已有既有账本来源；本分解节点的参与者、执行范围和时间先后仍须逐案核验。
\par

\noindent 在永嘉元年（145年），\eventanchor{EH-LEDGER-032-3}{冲帝夭亡，梁太后立八岁的刘缵为质帝，梁冀继续掌政}把东汉置于“冲帝夭亡，梁太后立八岁的刘缵为质帝，梁冀继续掌政”发生的政治与区域条件共同作用与该节点改变后续资源、联盟或制度处置的条件之间。可据《后汉书·冲帝纪》；《后汉书·质帝纪》；《后汉书·梁冀传》追踪其大体经过；核心事项已有既有账本来源；本分解节点的参与者、执行范围和时间先后仍须逐案核验。
\par

\noindent 东汉在永嘉元年（145年）的\eventanchor{EH-LEDGER-032-4}{“冲帝夭亡，梁太后立八岁的刘缵为质帝，梁冀继续掌政”的执行中的空间与人员处置}，并非孤立一事：“冲帝夭亡，梁太后立八岁的刘缵为质帝，梁冀继续掌政”发生的政治与区域条件共同作用，且该节点改变后续资源、联盟或制度处置的条件。材料见《后汉书·冲帝纪》；《后汉书·质帝纪》；《后汉书·梁冀传》；核心事项已有既有账本来源；本分解节点的参与者、执行范围和时间先后仍须逐案核验。
\par

\noindent \eventanchor{EH-LEDGER-032-5}{“冲帝夭亡，梁太后立八岁的刘缵为质帝，梁冀继续掌政”的后续制度或军政调整}发生于永嘉元年（145年），行动者为东汉。它回应核心节点发生后仍须处理其遗留问题，并把后续影响须在不同地区和材料类型中分别核验留给随后的人群和官署。《后汉书·冲帝纪》；《后汉书·质帝纪》；《后汉书·梁冀传》可作核验，惟核心事项已有既有账本来源；本分解节点的参与者、执行范围和时间先后仍须逐案核验。
\par

\noindent 从核心节点发生后仍须处理其遗留问题出发，东汉于永嘉元年（145年）出现\eventanchor{EH-LEDGER-032-6}{“冲帝夭亡，梁太后立八岁的刘缵为质帝，梁冀继续掌政”的异源材料与影响范围的复核}；后续影响须在不同地区和材料类型中分别核验。这里采用《后汉书·冲帝纪》；《后汉书·质帝纪》；《后汉书·梁冀传》的记载范围，不越过核心事项已有既有账本来源；本分解节点的参与者、执行范围和时间先后仍须逐案核验。
\par

\noindent 本初元年（146年），东汉的\eventanchor{EH-LEDGER-033-1}{“质帝被毒杀后，梁太后立刘志为桓帝，梁冀继续控制朝政”的前期动员与资源准备}使当时的权力、资源与信息约束推动相关过程展开转为可见的行动，这一过程为“质帝被毒杀后，梁太后立刘志为桓帝，梁冀继续控制朝政”提供可观察的条件。其依据为《后汉书·质帝纪》；《后汉书·桓帝纪》；《后汉书·梁冀传》，核心事项已有既有账本来源；本分解节点的参与者、执行范围和时间先后仍须逐案核验。
\par

\noindent \eventanchor{EH-LEDGER-033-2}{“质帝被毒杀后，梁太后立刘志为桓帝，梁冀继续控制朝政”的命令、联盟或地方响应}这一本初元年（146年）的节点，关联东汉、当时的权力、资源与信息约束推动相关过程展开和这一过程为“质帝被毒杀后，梁太后立刘志为桓帝，梁冀继续控制朝政”提供可观察的条件。《后汉书·质帝纪》；《后汉书·桓帝纪》；《后汉书·梁冀传》记录了它的基本轮廓；核心事项已有既有账本来源；本分解节点的参与者、执行范围和时间先后仍须逐案核验。
\par

\noindent 东汉在本初元年（146年）处理“质帝被毒杀后，梁太后立刘志为桓帝，梁冀继续控制朝政”发生的政治与区域条件共同作用时，出现\eventanchor{EH-LEDGER-033-3}{质帝被毒杀后，梁太后立刘志为桓帝，梁冀继续控制朝政}；该节点改变后续资源、联盟或制度处置的条件。可核材料为《后汉书·质帝纪》；《后汉书·桓帝纪》；《后汉书·梁冀传》，且核心事项已有既有账本来源；本分解节点的参与者、执行范围和时间先后仍须逐案核验。
\par

\noindent 本初元年（146年）的\eventanchor{EH-LEDGER-033-4}{“质帝被毒杀后，梁太后立刘志为桓帝，梁冀继续控制朝政”的执行中的空间与人员处置}，让东汉的“质帝被毒杀后，梁太后立刘志为桓帝，梁冀继续控制朝政”发生的政治与区域条件共同作用在具体事务中显形；该节点改变后续资源、联盟或制度处置的条件。《后汉书·质帝纪》；《后汉书·桓帝纪》；《后汉书·梁冀传》所能支持的范围仍受核心事项已有既有账本来源；本分解节点的参与者、执行范围和时间先后仍须逐案核验限制。
\par

\noindent 对东汉而言，\eventanchor{EH-LEDGER-033-5}{“质帝被毒杀后，梁太后立刘志为桓帝，梁冀继续控制朝政”的后续制度或军政调整}发生在本初元年（146年）：核心节点发生后仍须处理其遗留问题，随后后续影响须在不同地区和材料类型中分别核验。参看《后汉书·质帝纪》；《后汉书·桓帝纪》；《后汉书·梁冀传》；核心事项已有既有账本来源；本分解节点的参与者、执行范围和时间先后仍须逐案核验。
\par

\noindent 核心节点发生后仍须处理其遗留问题之下，东汉于本初元年（146年）经历\eventanchor{EH-LEDGER-033-6}{“质帝被毒杀后，梁太后立刘志为桓帝，梁冀继续控制朝政”的异源材料与影响范围的复核}，并面对后续影响须在不同地区和材料类型中分别核验。此处依《后汉书·质帝纪》；《后汉书·桓帝纪》；《后汉书·梁冀传》，核心事项已有既有账本来源；本分解节点的参与者、执行范围和时间先后仍须逐案核验。
\par
